\documentclass{jsarticle}
\usepackage{ascmac,amsmath,amssymb,amsthm,bm,physics,arydshln,mathcomp}
\newtheorem*{answer*}{解答}
\begin{document}

\begin{flushright}
6122111 三森誠真
\end{flushright}

\begin{itembox}[l]{\textbf{問2.91}}
ベクトル空間の間の線形写像$f : U \to V$ と$U$ の部分空間$W$ に対し, $W \subset \ker f$ を満たすとする. この時, 写像$\bar{f}$ : $U/W \to V$ を
$\bar{f}$ (cl($\bm{u}$)) = $f(\bm{u}) \ (\bm{u} \in U)$ と定めると, $\bar{f}$ はwell-defined な線形写像になることを証明せよ.
\end{itembox}

\begin{answer*}
\end{answer*}
cl($\bm{x}$) = cl($\bm{y}$) \ ($\bm{x}, \bm{y} \in U$) とする. このとき, $\bm{x} - \bm{y} \in W$, 即ち$\bm{x} - \bm{y} \in \ker f$ なので $f(\bm{x} - \bm{y}) = \bm{0}$. $f$ は

線形写像なので, $f(\bm{x}) = f(\bm{y})$. よって$\bar{f}$ (cl($\bm{x}$)) = $\bar{f}$ (cl($\bm{y}$)) となり$\bar{f}$ の定義は代表元の取り方に依らない. 

したがって$\bar{f}$ はwell-defined.

さらに, $\bar{f}$ (cl($\bm{x}$) + cl($\bm{y}$)) = $\bar{f}$ (cl($\bm{x}$ + $\bm{y}$)) = $f(\bm{x} + \bm{y})$ = $f(\bm{x})$ + $f(\bm{y})$ = $\bar{f}$ (cl($\bm{x}$)) + $\bar{f}$ (cl($\bm{y}$)),

$\forall c \in K$, \ $\bar{f}$ ($c \cdot$ cl($\bm{x}$)) = $\bar{f}$ (cl($c\bm{x}$)) = $f(c\bm{x})$ = $cf(\bm{x})$ = $c\bar{f}$ (cl($\bm{x}$)) となるので$\bar{f}$ は線形写像.

\bigskip

\begin{itembox}[l]{\textbf{問2.92}}
問2.91 において, $W = \ker f$ とする. この時, 写像$\bar{f}$ : $U/\ker f \to V$ の値域を制限した写像$\bar{f} : U/\ker f \to f(U)$ は同型写像となることを証明せよ. 
\end{itembox}

\begin{answer*}
\end{answer*}
問2.91より$\bar{f}$ が線形写像であることが分かるので$\bar{f}$ が全単射であることを示せばよい.

\textbf{全射性} : $\forall f(\bm{u}) \in f(U) \ (\bm{u} \in U), \ \exists$ cl($\bm{u}) \in U/\ker f$ s.t. $\bar{f}$ (cl($\bm{u}$)) = $f(\bm{u})$ なので全射.

\textbf{単射性} : $\bar{f}$ は全射なので, $\bar{f}$ (cl($\bm{x}$)) = $\bm{0}$ \ (cl$(\bm{x}) \in U/\ker f$) とすると, 

$\hspace{4zw} f(\bm{x}) = \bm{0} \ \Longrightarrow \ \bm{x} - \bm{0} \in \ker f \  \Longrightarrow$ \ cl($\bm{x}$) = cl($\bm{0}$) \ $\Longrightarrow \ker \bar{f}$ = cl($\bm{0}$) なので単射. ($\because$ 問2.66)

\end{document}
